%%=============================================================================
%% Homografietransformatie
%%=============================================================================

\section{Homografietransformatie: van pixel naar meter}
\label{sec:homografie}

% TODO: Literatuurstudie over formaat van de trampoline
Nu elk videoframe vrij is van lensvervorming, worden de pixelcoördinaten $(u, v)$ uit het beeldvlak via een basistransformatie geprojecteerd naar fysieke coördinaten $(x, y)$ op het trampolinedoek, uitgedrukt in meters. Aangezien de afmetingen van een wedstrijdtrampoline gestandaardiseerd zijn, kan er gebruik worden gemaakt van een \textit{planaire homografietransformatie}.

Een homografie $H$ is een $3 \times 3$ matrix die de projectieve transformatie beschrijft tussen twee vlakken: het camerabeeld en het trampolinedoek.

De transformatie wordt bepaald door vier corresponderende puntenparen:
\begin{enumerate}
    \item De vier hoekpunten van de trampolinemat worden in het videoframe (pixelcoördinaten) bepaald aan de hand van een \textit{landmark detection} model. Dit elimineert de nood aan manuele kalibratie en maakt de aanpak robuuster voor bewegende camerabeelden.
    \item De re{\"e}le afmetingen van de trampoline (bijv. 4.28\,m $\times$ 2.14\,m) worden als doelcoördinaten opgegeven. % TODO: Concrete afmetingen van de trampoline
\end{enumerate}

Met behulp van \texttt{cv2.getPerspectiveTransform} of \texttt{cv2.findHomography} wordt de matrix $H$ berekend. Vervolgens kan elke pixel $(u, v)$ op het trampolinedoek worden omgezet naar een coördinaat $(x, y)$ via:

\[
\begin{pmatrix} x' \\ y' \\ w \end{pmatrix} = H \cdot \begin{pmatrix} u \\ v \\ 1 \end{pmatrix}, \quad (x, y) = \left(\frac{x'}{w},\, \frac{y'}{w}\right)
\]

De nauwkeurigheid van de homografie wordt ge{\"e}valueerd door meetpunten op bekende posities op de mat te projecteren en te vergelijken met de gecorrigeerde re{\"e}le afstand. De verwachte maximale afwijking bedraagt enkele centimeters, afhankelijk van de cameraresolutie en de nauwkeurigheid van de handmatige hoekselectie.

% TODO: Figuur die de homografie toont + plot van de evaluatie